% Macros and settings shared by the pdf and the web version.

% Theorem-like environments. Those listed in the dependency graph by default
% are definition, lemma, proposition, theorem and corollary. Remarks are
% deliberately *not* part of the graph: we use them for commentary and for
% "digest notes" recording questions about the manuscript.
\newtheorem{theorem}{Theorem}[chapter]
\newtheorem{proposition}[theorem]{Proposition}
\newtheorem{lemma}[theorem]{Lemma}
\newtheorem{corollary}[theorem]{Corollary}

\theoremstyle{definition}
\newtheorem{definition}[theorem]{Definition}

\theoremstyle{remark}
\newtheorem{remark}[theorem]{Remark}

% Number sets
\newcommand{\T}{\mathbb{T}}
\newcommand{\R}{\mathbb{R}}
\newcommand{\Z}{\mathbb{Z}}
\newcommand{\N}{\mathbb{N}}
\newcommand{\C}{\mathbb{C}}

% Operators and notation
\newcommand{\Div}{\operatorname{div}}
\newcommand{\sym}{\operatorname{sym}}
\newcommand{\Ran}{\operatorname{Ran}}
\newcommand{\Jac}{\operatorname{Jac}}
\newcommand{\sech}{\operatorname{sech}}
\newcommand{\arctanh}{\operatorname{arctanh}}
\newcommand{\FP}{\operatorname{FP}}
\newcommand{\Var}{\operatorname{Var}}
\newcommand{\ip}[2]{\langle #1,#2\rangle}
\newcommand{\Dom}{\mathcal{D}}
\newcommand{\Cell}{\mathcal{C}}
